\section{Computational Resources and Project Management}
\label{sec:app1_resources_management}

\subsection{Conception of Novel Solutions}
We started from the isotropy literature \citep{ethayarajh-2019-contextual,rudman-etal-2022-isoscore}, which studies English text with older encoders. Two gaps followed: no study covers the scripts and tokenization levels of Indian languages, and no study compares a contrastive embedding model with a generative model of the same family. Early runs showed a sharp change in every metric at the final norm, so we added the pre/post-norm split and the norm breakdown to the method.

\subsection{Computational Environment}
Table~\ref{tab:app1_compute_resources} lists the resources used. The study needs one GPU at a time. Several models can run in parallel on separate GPUs with \texttt{--device}.

\begin{table}[t]
\centering
\small
\begin{tabular}{l>{\raggedright\arraybackslash}p{3.6cm}}
\toprule
Resource Component & Specification \\
\midrule
Hardware Accelerator & [GPU model] \\
GPU Video Memory & [VRAM]; float32 weights need 1.2 to 4.0\,GB \\
System Memory (RAM) & [RAM] \\
CPU Cores & [cores] \\
Persistent Storage & About 130\,MB of results per model, plus model weights and dataset cache \\
Software Stack & Python, PyTorch 2.10 to 2.11, Transformers 5.16 to 5.18, sentence-transformers 5.4 to 5.7, Datasets 4.8 to 5.0 \\
\bottomrule
\end{tabular}
\caption{Computational resources and environment for Approach 1.}
\label{tab:app1_compute_resources}
\end{table}

\subsection{Task Distribution and Workflow}
% TODO: fill in who did what for Approach 1.
[Detail the breakdown of responsibilities and workflow distribution among the team members assigned to Approach 1.]
